\documentclass[11pt,letterpaper]{article}
\usepackage[T1]{fontenc}
\usepackage{lmodern,amsmath,amssymb,microtype}
\usepackage[margin=1in]{geometry}
\usepackage[hidelinks]{hyperref}
\title{An 81-point unit-distance fixture that separates additive and ordinary colorings}
\author{Aafi Khan Malek\\\small Independent Researcher, Ahmedabad, India\\
\small\href{mailto:aafimalek2032@gmail.com}{aafimalek2032@gmail.com}}
\date{November 2026}
\begin{document}
\maketitle
This finite unit-distance graph has chromatic number four. It
nevertheless admits no proper coloring obtained by restricting
an additive homomorphism into a group of at most five elements.
It thus tests an algebraic search criterion, but does not
supply a six-chromatic example or a new plane bound.

Write \(s=\sqrt{-3}\), \(\omega=(1+s)/2\),
\(W=\{0,1,\omega,\ldots,\omega^5\}\) and \(T=W+W\).
The 19-point hexagonal ball \(T\) lies in \(\mathbb Q(s)\).
Let
\[
 \sigma=\frac{5+\sqrt{-11}}6,\qquad
 \rho=\frac{7+\sqrt{-15}}8,\qquad
 \tau=\frac{49+3\sqrt{-11}}{50},\qquad K=\mathbb Q(s,\sqrt{33}).
\]
They have modulus one. The base \(B=T\cup\sigma T\) has 37
sites since \(\sigma\notin\mathbb Q(s)\), and
\(F=B\cup\rho B\) has 73 sites since \(\rho\notin K\).
In the exterior join, all six pairs \(2u,2\rho u\) for sixth
roots \(u\) are actual unit edges:
\[
 |2u-2\rho u|^2=4(2-2\cdot7/8)=1.
\]
This uses the reference-radius join explained in
\href{https://arxiv.org/html/2010.12665v2#S4}{Parts, arXiv:2010.12665v2, \S4}.
The final fixture is
\[
 F^+=F\cup\{3,4,5\}\cup\{\tau,2\tau,3\tau,4\tau,5\tau\}.
\]
All eight sites are new: \(\tau\) and \(\tau/\sigma\) do not
belong to \(\mathbb Q(s)\); nonzero points of \(\rho B\)
lie outside \(K\); the new real sites have radius above two.
There is a further unit edge between \(5\) and \(5\tau\) since
\(25(2-2\cdot49/50)=1\).

The two diamonds in \(B\) with endpoints \(0,s\) and \(0,\sigma s\)
and middle vertices \((s\pm1)/2,\sigma(s\pm1)/2\) form a Moser
spindle, because \(|s-\sigma s|^2=1\). Hence no ordinary
three-coloring exists. For an additive homomorphism into the
binary group of order four, the doubled unit displacement
\(2(1-\rho)\) maps to zero. In the cyclic group of order four,
the edges for \(u=1,\omega,\omega-1\) require the three
images of \((1-\rho)u\) to be odd, impossible because the last
is the difference of the first two. The group of order five
kills \(5(1-\tau)\). All the undoubled or undivided differences
used here are in the additive group generated by the sites.
This proves the asserted restriction on additive maps.

For ordinary coloring, the accompanying exact replay
constructs all 81 distinct coordinates in a multiquadratic
basis, checks all 3,240 unordered distances by two algebraic
formulas, and validates a complete four-coloring on the
201 unit edges. Together with the spindle, this proves
\(\chi(F^+)=4\). In particular giving any two chosen
nonadjacent sites a fresh fifth color makes them monochromatic,
so this fixture does not force apart a marked pair under
ordinary five-coloring. The replay needs no solver after
the witness has been obtained. The proof and verification
were developed with AI assistance; neither is a Lean proof.
\end{document}
